\documentclass[11pt]{article}

\usepackage[a4paper,margin=1in]{geometry}
\usepackage{fancyhdr}
\usepackage{lastpage}
\usepackage{titlesec}
\usepackage{booktabs} 
\usepackage{enumitem} 
\usepackage{hyperref}
\usepackage{graphicx}
\usepackage{amsmath}
\usepackage{xcolor}
\usepackage{listings}

% Dark Theme Colors
\definecolor{pageblack}{rgb}{0,0,0}
\definecolor{textgrey}{rgb}{0.5,0.5,0.5}
\definecolor{codegreen}{rgb}{0,0.6,0}
\definecolor{codegray}{rgb}{0.5,0.5,0.5}
\definecolor{codepurple}{rgb}{0.58,0,0.82}
\definecolor{backcolour}{rgb}{0.05,0.05,0.05} 

% Apply Dark Theme
\pagecolor{pageblack}
\color{textgrey}

% Setup listings for CUDA/C++
\lstset{
    language=C++,
    backgroundcolor=\color{backcolour},   
    commentstyle=\color{codegreen},
    keywordstyle=\color{magenta},
    numberstyle=\tiny\color{codegray},
    stringstyle=\color{codepurple},
    basicstyle=\ttfamily\footnotesize\color{textgrey},
    breakatwhitespace=false,         
    breaklines=true,                 
    captionpos=b,                    
    keepspaces=true,                 
    numbers=left,                    
    numbersep=5pt,                  
    showspaces=false,                
    showstringspaces=false,
    showtabs=false,                  
    tabsize=2
}

% Header and Footer Setup
\pagestyle{fancy}
\fancyhf{}
\lhead{\color{textgrey}Project 4}
\chead{\color{textgrey}CUDA Optimization}
\rhead{\color{textgrey}Semester: SPRING 2026}
\lfoot{\color{textgrey}Arika Khor}
\cfoot{}
\rfoot{\color{textgrey}Page \thepage\ of \pageref{LastPage}}

\begin{document}

\begin{center}
    {\Large \textbf{\color{white}Project 4 Report: Tree Reduction and Stencil Computation in GPU}}\\
    \vspace{0.5em}
    Author: Arika Khor
\end{center}

\vspace{1em}

\section*{\color{white}Introduction}
This project explores the transformation of compute-intensive serial algorithms into high-performance CUDA applications. We focus on two domains: parallel search (PDNcoin mining) and image processing (convolutional stencils). The implementations leverage tree reduction, shared memory tiling, and atomic synchronization to maximize GPU utilization.

\noindent\rule{\textwidth}{0.4pt}

\section*{\color{white}Problem 1 \& 2: PDNcoin Mining and Reduction}

\subsection*{Algorithm Description}
The goal of these problems is to optimize a cryptocurrency mining workflow using CUDA. The process is divided into three parallelizable stages:
\begin{enumerate}[label=\arabic*.]
    \item \textbf{Nonce Generation:} Embarrassingly parallel generation of random trial values.
    \item \textbf{Hash Calculation:} Compute-intensive modular arithmetic per nonce.
    \item \textbf{Global Minimum Search:} Finding the winning nonce with the lowest hash value.
\end{enumerate}

\subsection*{Comparison: Implementation Differences}
The fundamental difference between Problem 1 and Problem 2 lies in the \textbf{offloading of the reduction step}.
\begin{itemize}
    \item \textbf{Problem 1 (Hybrid GPU/CPU):} Steps 1 and 2 are offloaded to the GPU. However, the GPU must copy the entire array of calculated hashes back to the host. At 10 million trials, this involves moving 40 MB over the PCIe bus, followed by an $O(N)$ serial loop on the host to find the minimum.
    \item \textbf{Problem 2 (Full GPU):} All three steps are offloaded. We implemented a tree reduction kernel that consolidated millions of results into a single 8-byte value directly in VRAM. This effectively eliminates the D2H (Device-to-Host) transfer bottleneck and the CPU search phase.
\end{itemize}

\subsection*{Optimized Reduction with Atomic Tie-Breaking}
To find the global minimum while ensuring the smallest nonce is picked in case of a hash tie, we utilized a packed 64-bit atomic approach.

\begin{lstlisting}
// Packing Hash (Upper 32) and Nonce (Lower 32) for atomic priority
if (tid == 0) {
    unsigned long long local_min = ((unsigned long long)s_hashes[0] << 32) | (unsigned long long)s_nonces[0];
    unsigned long long assumed;
    unsigned long long old = *d_min_result;
    do {
        assumed = old;
        if (local_min < assumed) {
            old = atomicCAS(d_min_result, assumed, local_min);
        } else break;
    } while (assumed != old);
}
\end{lstlisting}

\subsection*{Performance Results (Mining Trials)}
\begin{center}
\begin{tabular}{@{}lcc@{}}
\toprule
Implementation & 5 Million Trials & 10 Million Trials \\
\midrule
\textbf{Serial Mining (CPU)} & 588.61s & 1177.22s \\
\textbf{Problem 1 (GPU Hash)} & 0.431s & 0.777s \\
\textbf{Problem 2 (GPU Reduction)} & 0.440s & 0.754s \\
\bottomrule
\end{tabular}
\end{center}

\noindent\rule{\textwidth}{0.4pt}

\section*{\color{white}Problem 3 \& 4: Tiled Convolution and Max-Pooling}

\subsection*{Algorithm Description}
These problems implement a deep learning inference pipeline using 2D stencils. We implemented a 2D convolution layer followed by a max-pooling operation.

\subsection*{Comparison: Implementation Differences}
\begin{itemize}
    \item \textbf{Problem 3 (Spatial Optimization):} Focuses on optimizing memory bandwidth using **shared memory tiling**. By loading tiles collaboratively, threads reuse data for stencil calculations, drastically reducing redundant global memory requests.
    \item \textbf{Problem 4 (Integrated Pipeline):} Combines the convolution from Problem 3 with a max-pooling stage. The convolution output is kept in global memory and immediately read by the pooling kernel. This "chaining" avoids intermediate host interaction and keeps the full feature extraction on the GPU.
\end{itemize}

\subsection*{Collaborative Halo Loading (Problem 3)}
To handle the 5x5 filter, each thread block loads a "halo" of surrounding pixels. We used dynamic shared memory to support runtime-configurable block sizes and radii.

\begin{lstlisting}
// collaborative loading of tile + halo into dynamic shared memory
for (int i = ty; i < tile_dim; i += blockDim.y) {
    for (int j = tx; j < tile_dim; j += blockDim.x) {
        int global_row = blockIdx.y * blockDim.y + i;
        int global_col = blockIdx.x * blockDim.x + j;
        if (global_row < n_row + 2 * radius && global_col < n_col + 2 * radius)
            tile_s[i * tile_dim + j] = paddedInput[global_row * (n_col + 2 * radius) + global_col];
        else tile_s[i * tile_dim + j] = 0;
    }
}
\end{lstlisting}

\subsection*{Max-Pooling Window Logic (Problem 4)}
The max-pooling kernel enhances feature clarity by searching for local maxima.

\begin{lstlisting}
if (row < n_row && col < n_col) {
    int max_val = INT_MIN;
    for (int i = -radius; i <= radius; i++) {
        for (int j = -radius; j <= radius; j++) {
            int cur_row = row + i;
            int cur_col = col + j;
            if (cur_row >= 0 && cur_row < n_row && cur_col >= 0 && cur_col < n_col) {
                int val = convolved_input[cur_row * n_col + cur_col];
                if (val > max_val) max_val = val;
            }
        }
    }
    output[row * n_col + col] = max_val;
}
\end{lstlisting}

\subsection*{Performance Results (2048 x 2048 Matrix)}
\begin{center}
\begin{tabular}{@{}lc@{}}
\toprule
Implementation Stage & Runtime (seconds) \\
\midrule
\textbf{Problem 3 (Convolution Only)} & 0.000671 \\
\textbf{Problem 4 (Integrated Pipeline)} & 0.000981 \\
\bottomrule
\end{tabular}
\end{center}

\noindent\rule{\textwidth}{0.4pt}

\section*{\color{white}AI Usage and Software Engineering Workflow}

\subsection*{Gemini CLI: The Interactive Engine}
Development was driven by the \textbf{Gemini CLI}, an agentic environment capable of direct workspace interaction, shell execution, and surgical code refactoring.

\subsection*{The GSD (Get Shit Done) Framework}
To manage the complexity of this project, I used the **GSD (Get Shit Done)** workflow:
\begin{enumerate}[label=\alph*.]
    \item \textbf{Initialization:} Automatic instruction mapping to \texttt{ROADMAP.md} and \texttt{REQUIREMENTS.md}.
    \item \textbf{Research:} Bottleneck analysis (identifying the D2H PCIe overhead).
    \item \textbf{Planning:} Decomposing problems into atomic, verifiable phases.
    \item \textbf{Execution:} Wave-based implementation with GTest verification checkpoints.
\end{enumerate}

\noindent\rule{\textwidth}{0.4pt}

\section*{\color{white}Project Layout}
\begin{itemize}
    \item \textbf{common/}: Shared `gpuErrchk` error logic and timing functions.
    \item \textbf{Problem\_1/2}: Mining logic comparison (Fat vs. Thin D2H transfers).
    \item \textbf{Problem\_3/4}: Image processing (Spatial tiling vs. Chained pipelines).
    \item \textbf{serial/}: Centralized reference C programs for performance baselines.
    \item \textbf{tests/}: Comprehensive Google Test suite for all kernels.
\end{itemize}

\end{document}
